\section{Conclusion}
\label{sec:conclusion}

We presented a multi-agent simulation framework where LLM-driven provider and insurer agents negotiate prior authorization for a specialty medication under information asymmetry. In a 3$\times$3 factorial experiment on a single gray-zone clinical case, we found that the insurer's choice between requesting additional information and outright denial at initial review is the mechanism through which gray-zone ambiguity is resolved. Default (no-guidance) insurer agents denied the drug while approving diagnostics, applying policy literally rather than interpretively, suggesting that AI adjudication systems deployed without explicit guidance may default to denial in ambiguous cases. Game-theoretic analysis shows the game's classification is parameter-dependent: at observed administrative costs (\$5.79/turn for providers vs.\ \$0.05/turn for insurers, a 116$\times$ asymmetry), the game is an asymmetric exploitation where denial is nearly costless for insurers; under richer cost models incorporating litigation, reputational, and network costs, it can transition to a prisoner's dilemma. Extending this framework to multiple clinical cases, policy types, foundation models, and replicated runs would characterize the generality of these findings.